
\documentclass{article}
\usepackage{colortbl}
\usepackage{makecell}
\usepackage{multirow}
\usepackage{supertabular}

\begin{document}

\newcounter{utterance}

\centering \large Interaction Transcript for game `wordle', experiment `high\_frequency\_words\_no\_clue\_no\_critic', episode 8 with Qwen3.6{-}35B{-}A3B{-}FP8{-}without{-}reasoning.
\vspace{24pt}

{ \footnotesize  \setcounter{utterance}{1}
\setlength{\tabcolsep}{0pt}
\begin{supertabular}{c@{$\;$}|p{.15\linewidth}@{}p{.15\linewidth}p{.15\linewidth}p{.15\linewidth}p{.15\linewidth}p{.15\linewidth}}
    \# & \multicolumn{2}{c}{Player} && \multicolumn{2}{c}{Game Master} \\
    \hline

    \theutterance \stepcounter{utterance}  
    & & & \multicolumn{4}{p{0.6\linewidth}}{
        \cellcolor[rgb]{0.9,0.9,0.9}{
            \makecell[{{p{\linewidth}}}]{
                \texttt{\tiny{[P1$\langle$GM]}}
                \texttt{Verilen kuralları kullanarak kelime tahmin etmeyi seven bir dil sihirbazısın.} \\
\\ 
\texttt{Wordle’a hoş geldin! Hedef kelimeyi tahmin etmek için altı denemen var: beş küçük harften oluşan, geçerli bir kelime. Her tahmin için kısa bir açıklama sunmak üzere lütfen "açıklama:" ve "tahmin:" etiketlerini kullan.} \\
\\ 
\texttt{Örneğin, tahminin "yemek" ise, yanıtın şöyle olmalı:} \\
\texttt{açıklama: bu yaygın bir beş harfli kelime ve tahminime bu kelimeyle başlıyorum.} \\
\texttt{tahmin: yemek} \\
\\ 
\texttt{Her tahminden sonra yanıtın doğrulanacak ve hangi harflerin doğru olduğunu (yeşil), hangi harflerin doğru ama yanlış konumda olduğunu (sarı) ve hangi harflerin yanlış olduğunu (kırmızı) belirten geri bildirim alacaksın. Bu geri bildirim, bir sonraki tahmininde hangi harfleri dahil edeceğini veya hariç tutacağını belirlemede faydalı olabilir.} \\
\\ 
\texttt{Örneğin, "yemek" için geri bildirim şöyle olabilir:} \\
\texttt{tahmin\_geri\_bildirimi: y<sarı> e<sarı> m<yeşil> e<sarı> k<kırmızı>} \\
\\ 
\texttt{Açıklama, yeni bir tahmine ulaşmak için tahmin\_geri\_bildirimi’in nasıl kullanıldığına dair ayrıntılar içermelidir.} \\
\\ 
\texttt{Başka hiçbir metin üretme.} \\
\\ 
\texttt{İlk tahmininle başlayalım.} \\
            }
        }
    }
     \\ \\

    \theutterance \stepcounter{utterance}  
    & \multicolumn{4}{p{0.6\linewidth}}{
        \cellcolor[rgb]{0.561,0.737,0.561}{
            \makecell[{{p{\linewidth}}}]{
                \texttt{\tiny{[P1$\rangle$GM]}}
                \texttt{açıklama: Bu yaygın bir beş harfli kelime ve tahminime bu kelimeyle başlıyorum.} \\
\texttt{tahmin: araba} \\
            }
        }
    }
    & & \\ \\

    \theutterance \stepcounter{utterance}  
    & & & \multicolumn{2}{p{0.3\linewidth}}{
        \cellcolor[rgb]{0.9,0.9,0.9}{
            \makecell[{{p{\linewidth}}}]{
                \texttt{\tiny{[GM$|$GM]}}
                \texttt{attempts = 1} \\
\texttt{target\_word = şehir} \\
\texttt{guess = araba} \\
\texttt{guess\_feedback = a<red> r<yellow> a<red> b<red> a<red>} \\
            }
        }
    }
    & & \\ \\

    \theutterance \stepcounter{utterance}  
    & & & \multicolumn{4}{p{0.6\linewidth}}{
        \cellcolor[rgb]{0.9,0.9,0.9}{
            \makecell[{{p{\linewidth}}}]{
                \texttt{\tiny{[P1$\langle$GM]}}
                \texttt{tahmin\_geri\_bildirimi: a<kırmızı> r<sarı> a<kırmızı> b<kırmızı> a<kırmızı>} \\
\\ 
\texttt{Yanıtını yalnızca bu biçimde ver.} \\
\texttt{açıklama: ayrıntılar} \\
\texttt{tahmin: kelime} \\
            }
        }
    }
     \\ \\

    \theutterance \stepcounter{utterance}  
    & \multicolumn{4}{p{0.6\linewidth}}{
        \cellcolor[rgb]{0.561,0.737,0.561}{
            \makecell[{{p{\linewidth}}}]{
                \texttt{\tiny{[P1$\rangle$GM]}}
                \texttt{açıklama: İlk tahminde 'a' harfi tamamen elendi ve 'r' harfi yanlış konumda. 'b' harfi de elendi. 'r' harfini içeren ve 'r' harfinin ikinci veya üçüncü harf olabileceği, 'a' ve 'b' içermeyen yaygın bir kelime deneyeceğim. "kiraz" kelimesi 'r' harfini içerir ve 'a' harfini içermez.} \\
\texttt{tahmin: kiraz} \\
            }
        }
    }
    & & \\ \\

    \theutterance \stepcounter{utterance}  
    & & & \multicolumn{2}{p{0.3\linewidth}}{
        \cellcolor[rgb]{0.9,0.9,0.9}{
            \makecell[{{p{\linewidth}}}]{
                \texttt{\tiny{[GM$|$GM]}}
                \texttt{attempts = 2} \\
\texttt{target\_word = şehir} \\
\texttt{guess = kiraz} \\
\texttt{guess\_feedback = k<red> i<yellow> r<yellow> a<red> z<red>} \\
            }
        }
    }
    & & \\ \\

    \theutterance \stepcounter{utterance}  
    & & & \multicolumn{4}{p{0.6\linewidth}}{
        \cellcolor[rgb]{0.9,0.9,0.9}{
            \makecell[{{p{\linewidth}}}]{
                \texttt{\tiny{[P1$\langle$GM]}}
                \texttt{tahmin\_geri\_bildirimi: k<kırmızı> i<sarı> r<sarı> a<kırmızı> z<kırmızı>} \\
\\ 
\texttt{Yanıtını yalnızca bu biçimde ver.} \\
\texttt{açıklama: ayrıntılar} \\
\texttt{tahmin: kelime} \\
            }
        }
    }
     \\ \\

    \theutterance \stepcounter{utterance}  
    & \multicolumn{4}{p{0.6\linewidth}}{
        \cellcolor[rgb]{0.561,0.737,0.561}{
            \makecell[{{p{\linewidth}}}]{
                \texttt{\tiny{[P1$\rangle$GM]}}
                \texttt{açıklama: İlk tahminde 'r' harfi ikinci konumdaydı (sarı). İkinci tahminde 'r' üçüncü konumdaydı (sarı) ve 'i' ikinci konumdaydı (sarı). Bu, 'r' harfinin kelimenin ilk veya son konumunda olabileceğini, 'i' harfinin ise ilk veya son konumda olabileceğini gösteriyor. 'a', 'b', 'k', 'z' harfleri elendi. 'r' ve 'i' harflerini koruyarak, 'r'yi ilk harf, 'i'yi son harf olarak deneyen yaygın bir kelime "kirpi" olabilir ama 'p' ve 'y' yeni harfler. Daha yaygın bir kelime arıyorum. 'r' ilk harf, 'i' son harf olabilir: "r...i". 'r' ikinci harf, 'i' ilk harf olabilir: "i...r". 'r' son harf, 'i' ilk harf olabilir: "i..r.". 'r' son harf, 'i' son harf olamaz. 'r' ilk harf, 'i' ikinci harf olabilir: "ri...". "ri...i" olamaz çünkü 'i' sadece bir kez geçiyor gibi görünüyor (genelde Wordle'da harf sayısı sınırlıdır, ama burada 'i' sarı olduğu için var). 'r' ikinci harf, 'i' üçüncü harf olabilir. "i" sarı, "r" sarı. "kiraz"da 'i' 2. 'r' 3. idi. "araba"da 'r' 2. idi. 'r' 1. veya 5. olabilir. 'i' 1. veya 5. olabilir. "r" 1. ise "i" 5. olabilir: "r...i". "r" 5. ise "i" 1. olabilir: "i..r". "r" 1. ise "i" 1. olamaz. "r" 5. ise "i" 5. olamaz. "r" 2. ise "i" 1. veya 5. olabilir. "} \\
            }
        }
    }
    & & \\ \\

    \theutterance \stepcounter{utterance}  
    & & & \multicolumn{2}{p{0.3\linewidth}}{
        \cellcolor[rgb]{0.9,0.9,0.9}{
            \makecell[{{p{\linewidth}}}]{
                \texttt{\tiny{[GM$|$GM]}}
                \texttt{Error: The response should contain the 'tahmin:' keyword exactly once.} \\
            }
        }
    }
    & & \\ \\

    \theutterance \stepcounter{utterance}  
    & & & \multicolumn{2}{p{0.3\linewidth}}{
        \cellcolor[rgb]{0.9,0.9,0.9}{
            \makecell[{{p{\linewidth}}}]{
                \texttt{\tiny{[GM$|$GM]}}
                \texttt{game\_result = ABORT} \\
            }
        }
    }
    & & \\ \\

\end{supertabular}
}

\end{document}
